% !TeX program = lualatex
% ================================================================
% Urdu Parallel Text Version of RecurringDecimalsToFractions
%
% Generated by the server using the following settings:
%
%   language            Urdu (urd)
%   text direction      right to left
%   font                Noto Nastaliq Urdu
%   fallback font       Noto Serif
%   built from          latex/worksheets/fractions/recurringDecimalsToFractions.tex
%   vocabulary key      latex/worksheets/fractions/recurringDecimalsToFractions/L2/urd/recurringDecimalsToFractions_urduKey.tex
%   tier 3 vocabulary   green   RGB 0 128 0
%   English text        black   RGB 0 0 0
%   Urdu text           purple  RGB 112 48 160
%   layout macros       version 3
%
% Please don't edit any information in this comment block - the server
% compares it against the current settings and rebuilds this file when
% they differ, so changes here would be overwritten. However, feel free
% to make updates to the LaTeX below if you want to edit it.
% ================================================================

\documentclass[12pt,a4paper]{article}
% ================================================================
% Copied from the English sheet this was translated from, so that
% anything it sets up for itself still works here
% ================================================================
\usepackage[a4paper,top=12mm,bottom=14mm,left=15mm,right=15mm]{geometry}
\usepackage[T1]{fontenc}
\usepackage{amsmath}
\usepackage{amssymb}
\usepackage{enumitem}
\usepackage{multicol}
\usepackage{tikz}
\usepackage[most]{tcolorbox}
\usetikzlibrary{decorations.pathreplacing,arrows.meta}

% ================================================================
% Recurring decimals to fractions.
%
% An extension sheet for the top end of a middle ability Year 9
% class, designed to be worked through alone: every method is
% introduced by a worked example, practised once with the steps
% written out, and then practised without them. Answers are on the
% last page so that pupils can check each part before going on.
%
% The sheet is written on, so it is set with open leading and wide
% list separation, and every question is followed by a working area
% sized to fill the rest of its page. The \newpage commands in the
% body are placed so that no question is ever split across the fold;
% because of them a part may run to more than one page.
% ================================================================

\pagestyle{empty}
\linespread{1.2}\selectfont
\setlength{\parskip}{5pt}
\renewcommand{\arraystretch}{1.35}
\setlist[enumerate,1]{label=\textbf{Q\arabic*.}, leftmargin=*, itemsep=12pt, topsep=8pt, parsep=4pt}
\setlist[enumerate,2]{label=(\alph*), leftmargin=*, itemsep=8pt, topsep=5pt, parsep=3pt}
\setlength{\multicolsep}{8pt}
\setlength{\parindent}{0pt}

\newcommand{\ansline}[1][2.2cm]{\underline{\hspace{#1}}}

% A worked example: the method shown in full before it is practised.
\newtcolorbox{wex}[1]{colback=blue!4,colframe=blue!45!black,boxrule=0.6pt,
  left=6pt,right=6pt,top=4pt,bottom=5pt,
  title={\bfseries\small Worked example: #1},fonttitle=\small,
  toptitle=1pt,bottomtitle=1pt}

% The idea a section is built on, stated once and boxed.
\newtcolorbox{keyidea}{colback=orange!7,colframe=orange!60!black,boxrule=0.6pt,
  left=6pt,right=6pt,top=4pt,bottom=4pt}

% The section headings, kept tight so the parts do not eat the page.
\newcommand{\heading}[1]{%
  \vspace{2pt}
  {\color{blue!45!black}\rule{\linewidth}{1pt}}\\[1pt]
  {\large\bfseries #1}\\[-6pt]
  {\color{blue!45!black}\rule{\linewidth}{0.4pt}}
  \vspace{2pt}}
\newcommand{\parttitle}[2]{\heading{Part #1: #2}}

% ----------------------------------------------------------------
% One line of a subtraction diagram. Every row puts its decimal
% point at the same x, so that recurring tails which really are
% identical end up in the same columns on the page.
% \drow{<y>}{<label>}{<digits before the point>}{<digits after>}
\newcommand{\drow}[4]{%
  \node[anchor=east,font=\small] at (-1.25,#1) {#2};
  \node[anchor=east,font=\small] at (-0.12,#1) {#3};
  \node[font=\small] at (0,#1) {.};
  \foreach \d [count=\i from 0] in {#4}
    {\node[font=\small] at ({0.34+\i*0.34},#1) {\d};}%
}
% The band of columns from digit #1 to digit #2, used to shade a tail.
\newcommand{\tailbox}[4]{%
  ({0.34+#1*0.34-0.17},#3) rectangle ({0.34+#2*0.34+0.17},#4)}
% The same row as \drow, but with a writing line in each digit column so that
% a pupil fills the row in. #4 is how many columns to draw before the dots.
% \drowgap{<y>}{<label>}{<before the point>}{<number of columns>}
\newcommand{\drowgap}[4]{%
  \node[anchor=east,font=\small] at (-1.25,#1) {#2};
  \node[anchor=east,font=\small] at (-0.12,#1) {#3};
  \node[font=\small] at (0,#1) {.};
  \foreach \i in {1,...,#4}
    {\draw[gray!60] ({0.34+(\i-1)*0.34-0.13},{#1-0.16}) --
                    ({0.34+(\i-1)*0.34+0.13},{#1-0.16});}
  \node[font=\small] at ({0.34+#4*0.34},#1) {$\ldots$};%
}
% Solving by dividing both sides: the equation before above the equation
% after, with a curved arrow down each side carrying the division. Both
% arrows bend outwards, so neither crosses the equation.
% \divideboth{<lhs>}{<rhs>}{<new lhs>}{<new rhs>}{<divisor>}
\newcommand{\divideboth}[5]{%
  \begin{tikzpicture}
    \node (tl) at (0,0)     {$#1$};
    \node       at (0.7,0)  {$=$};
    \node (tr) at (1.5,0)   {$#2$};
    \node (bl) at (0,-1.5)  {$#3$};
    \node       at (0.7,-1.5) {$=$};
    \node (br) at (1.5,-1.5) {$#4$};
    \draw[-{Stealth[length=2mm]},blue!45!black] (tl) to[bend right=65]
      node[left=1pt,font=\scriptsize]{$\div\,#5$} (bl);
    \draw[-{Stealth[length=2mm]},blue!45!black] (tr) to[bend left=65]
      node[right=1pt,font=\scriptsize]{$\div\,#5$} (br);
  \end{tikzpicture}%
}

% Characters this language's font has not got are borrowed from another,
% rather than being silently dropped - see docs/translated-sheet-latex.md
\directlua{%
  if luaotfload and luaotfload.add_fallback then
    luaotfload.add_fallback("eallatinfallback",
      { "Noto Serif:mode=harf;" })
  else
    tex.error("this luaotfload is too old to register a fallback font")
  end
}
\usepackage[english,bidi=basic]{babel}
\babelprovide[import]{urdu}
\babelfont[urdu]{rm}[Renderer=Harfbuzz, BoldFont={Noto Nastaliq Urdu}, ItalicFont={Noto Nastaliq Urdu}, BoldItalicFont={Noto Nastaliq Urdu}, RawFeature={fallback=eallatinfallback}]{Noto Nastaliq Urdu}
\babelfont[urdu]{sf}[Renderer=Harfbuzz, BoldFont={Noto Nastaliq Urdu}, ItalicFont={Noto Nastaliq Urdu}, BoldItalicFont={Noto Nastaliq Urdu}, RawFeature={fallback=eallatinfallback}]{Noto Nastaliq Urdu}
\newcommand{\ealtext}[1]{\foreignlanguage{urdu}{#1}}
\newcommand{\ealtextblock}[1]{{\raggedleft\foreignlanguage{urdu}{#1}\par}}
% ================================================================
% EAL layout helpers
% ================================================================
\usepackage{xcolor}

\definecolor{ealkeycolour}{RGB}{0,128,0}
\definecolor{ealenglish}{RGB}{0,0,0}
\definecolor{ealtwocolour}{RGB}{112,48,160}

\newsavebox{\ealboxen}
\newsavebox{\ealboxtr}
\newlength{\ealglosswd}
\newlength{\ealglossraise}
\setlength{\ealglossraise}{1.15em}

% a tier 3 word, in English and in translation alike
\newcommand{\ealkey}[1]{\textcolor{ealkeycolour}{#1}}
\newcommand{\ealkeytr}[1]{\textcolor{ealkeycolour}{\ealtext{#1}}}

% One translated word above one English word. Built from kernel
% primitives, never a tabular, which would break wherever the sheet
% loads array, booktabs or tabularx. The box is as wide as the wider of
% the two, so a long translation pushes its neighbours aside instead of
% overlapping them, and it claims its full height so a table row or a
% TikZ node makes room for it.
\newcommand{\ealgloss}[2]{%
  \sbox{\ealboxen}{\ealkey{#1}}%
  \sbox{\ealboxtr}{\small\textcolor{ealtwocolour}{\ealtext{#2}}}%
  \setlength{\ealglosswd}{\wd\ealboxen}%
  \ifdim\wd\ealboxtr>\ealglosswd\setlength{\ealglosswd}{\wd\ealboxtr}\fi
  \leavevmode
  \makebox[\ealglosswd][c]{%
    \makebox[0pt][c]{\usebox{\ealboxen}}%
    \makebox[0pt][c]{\raisebox{\ealglossraise}{\usebox{\ealboxtr}}}%
  }%
}

% opens up line spacing around a block of glosses, so the raised
% translations sit clear of the line above
\newenvironment{ealglossed}{\par\linespread{2.1}\selectfont\ignorespaces}{\par}

% a whole translated sentence above its English counterpart. block level,
% so it does not belong inside a tabular cell
\newcommand{\ealpara}[2]{%
  \par\addvspace{0.5\baselineskip}%
  {\color{ealtwocolour}\ealtextblock{#1}}%
  \penalty200
  {\color{ealenglish}#2\par}%
  \addvspace{0.5\baselineskip}%
}

\begin{document}

\begin{center}
  {\Large\bfseries \ealkey{Recurring Decimals} into \ealkey{Fractions}}\\[8pt]
\end{center}
\vspace{1pt}

\begin{keyidea}
\small
\ealpara{ہر \ealkeytr{تکراری اعشاریہ} --- چاہے کتنا ہی طویل اور پیچیدہ ہو --- ایک \ealkeytr{کسر} کے طور پر لکھا جا سکتا ہے۔ اس شیٹ کی مثالوں کی پیروی کریں تاکہ \ealkeytr{تکراری اعشاریوں} کو \ealkeytr{کسروں} میں بدلنا سیکھیں۔}{Every \ealkey{recurring decimal} --- however long and however complicated --- can be written as a \ealkey{fraction}. Follow the examples in this sheet to learn how to convert \ealkey{recurring decimals} to \ealkey{fractions}.}
\end{keyidea}

\parttitle{1}{Reading and writing \ealkey{recurring decimals}}

\ealpara{جب آپ ایک پورے عدد کو دوسرے سے تقسیم کرتے ہیں، تو نقطے کے بعد والے \ealkeytr{ہندسے} یا تو رک جاتے ہیں، یا ایک ایسے گروہ میں گرتے ہیں جو ہمیشہ دہرتا رہتا ہے۔ دہرائے جانے والے گروہوں کو \textbf{پہلے} \ealkeytr{ہندسے} کے اوپر ایک نقطہ اور \textbf{آخری} \ealkeytr{ہندسے} کے اوپر ایک نقطہ لگا کر ظاہر کیا جاتا ہے۔}{When you divide one whole number by another, the \ealkey{digits} after the point either stop, or fall into a block that repeats forever. Repeating blocks are marked with a dot over the \textbf{first} \ealkey{digit} and a dot over the \textbf{last} \ealkey{digit}.}

\begin{center}
\begin{tikzpicture}[font=\small]
  \def\dgw{0.44}
  \foreach \s/\c in {0/blue!22, 2/blue!9, 4/blue!22, 6/blue!9}{
    \fill[\c] ({0.7+\s*\dgw-0.5*\dgw},-0.28) rectangle ({0.7+(\s+2)*\dgw-0.5*\dgw},0.3);}
  \node at (0,0) {$0$};
  \node at (0.3,0) {.};
  \foreach \d [count=\i from 0] in {2,7,2,7,2,7,2,7}
    {\node at ({0.7+\i*\dgw},0) {\d};}
  \node at ({0.7+8*\dgw+0.2},0) {$\ldots$};
  \draw[decorate,decoration={brace,amplitude=3pt}]
    ({0.7-0.5*\dgw},0.36) -- ({0.7+1.5*\dgw},0.36)
    node[midway,above=2pt,font=\scriptsize]{the block that repeats};
  \node[anchor=west,font=\small] at ({0.7+8*\dgw+0.75},0)
    {so we write it $0.\dot{2}\dot{7}$};
\end{tikzpicture}
\end{center}
\vspace{-4pt}

\ealpara{تو $0.\dot{6}=0.6666666\ldots$، \; $0.\dot{2}\dot{7}=0.2727272\ldots$، \;
$0.\dot{4}0\dot{9}=0.4094094\ldots$ \; اور $0.1\dot{6}=0.1666666\ldots$ ---
غور کریں کہ آخری میں صرف $6$ پر نقطہ ہے، اس لیے صرف $6$ دہرایا جاتا ہے۔ $1$ صرف دہائیوں کے کالم میں پایا جاتا ہے، اور کبھی واپس نہیں آتا۔}{So $0.\dot{6}=0.6666666\ldots$, \; $0.\dot{2}\dot{7}=0.2727272\ldots$, \;
$0.\dot{4}0\dot{9}=0.4094094\ldots$ \; and $0.1\dot{6}=0.1666666\ldots$ ---
notice that in the last one only the $6$ carries a dot, so only the $6$
repeats. The $1$ only is found in the tenths column, and never comes back.}

\vspace{2em}

\begin{enumerate}

\item \ealpara{ہر عدد کے پہلے آٹھ \ealkeytr{اعشاریہ مقامات} لکھیں۔}{Write out the first eight \ealkey{decimal places} of each number.}
\begin{multicols}{3}
\begin{enumerate}
  \item $0.\dot{5}=0.\ansline[2.3cm]$
  \item $0.\dot{1}\dot{8}=0.\ansline[2.3cm]$
  \item $0.3\dot{7}=0.\ansline[2.3cm]$
  \item $0.\dot{2}0\dot{6}=0.\ansline[2.3cm]$
  \item $0.4\dot{5}\dot{1}=0.\ansline[2.3cm]$
  \item $0.\dot{9}=0.\ansline[2.3cm]$
\end{enumerate}
\end{multicols}

\vspace{1em}

\item \ealpara{ان میں سے ہر ایک کو \ealkeytr{نقطہ علامت} استعمال کرتے ہوئے لکھیں۔}{Write each of these using \ealkey{dot notation}.}
\begin{multicols}{3}
\begin{enumerate}
  \item $0.8888\ldots=\ansline[1.5cm]$
  \item $0.636363\ldots=\ansline[1.5cm]$
  \item $0.24444\ldots=\ansline[1.5cm]$
  \item $0.157157\ldots=\ansline[1.5cm]$
  \item $0.91666\ldots=\ansline[1.5cm]$
  \item $0.40505\ldots=\ansline[1.5cm]$
\end{enumerate}
\end{multicols}

\vspace{1em}

\item
\begin{enumerate}
  \item \ealpara{کون بڑا ہے، $0.\dot{4}\dot{5}$ یا $0.4\dot{5}$؟ فیصلہ کرنے میں مدد کے لیے دونوں کو چھ \ealkeytr{اعشاریہ مقامات} تک لکھیں، اور اپنا جواب سمجھائیں۔}{Which is bigger, $0.\dot{4}\dot{5}$ or $0.4\dot{5}$? Write both out to six \ealkey{decimal places} to help you decide, and explain your answer.}

        \vspace{7em}
  \item \ealpara{ایک \ealkeytr{تکراری اعشاریہ} لکھیں جو $0.\dot{3}$ اور $0.\dot{4}$ کے درمیان ہو۔}{Write down a \ealkey{recurring decimal} that lies between $0.\dot{3}$ and $0.\dot{4}$.} \ansline[3cm]
\end{enumerate}

\end{enumerate}

% \parttitle{2}{The family of nines}

% Before any algebra, here is a pattern worth knowing. Use short division to
% check the first one or two, then follow the pattern.

% \begin{enumerate}[resume]

% \item Complete the following.
% \begin{multicols}{4}
% \begin{enumerate}
%   \item $\tfrac{1}{9}=0.\dot{1}$
%   \item $\tfrac{2}{9}=\ansline[1.4cm]$
%   \item $\tfrac{4}{9}=\ansline[1.4cm]$
%   \item $\tfrac{7}{9}=\ansline[1.4cm]$
% \end{enumerate}
% \end{multicols}

% \item Now with two and three nines on the bottom.
% \begin{multicols}{3}
% \begin{enumerate}
%   \item $\tfrac{1}{99}=0.\dot{0}\dot{1}$
%   \item $\tfrac{7}{99}=\ansline[1.8cm]$
%   \item $\tfrac{14}{99}=\ansline[1.8cm]$
%   \item $\tfrac{58}{99}=\ansline[1.8cm]$
%   \item $\tfrac{1}{999}=0.\dot{0}0\dot{1}$
%   \item $\tfrac{271}{999}=\ansline[1.8cm]$
% \end{enumerate}
% \end{multicols}

% \end{enumerate}

% \begin{keyidea}
% \small
% \textbf{The nines rule.} If a block of $n$ digits repeats \emph{straight after
% the decimal point}, the fraction is that block written over $n$ nines. So
% $0.\dot{8}=\tfrac{8}{9}$, \; $0.\dot{2}\dot{7}=\tfrac{27}{99}$ \; and \;
% $0.\dot{4}0\dot{9}=\tfrac{409}{999}$. Always simplify afterwards ---
% $\tfrac{27}{99}=\tfrac{3}{11}$.
% \end{keyidea}

% \begin{enumerate}[resume]

% \item Use the nines rule. Give each answer as a fraction in its simplest form.
% \begin{multicols}{3}
% \begin{enumerate}
%   \item $0.\dot{7}$ \ansline[1.6cm]
%   \item $0.\dot{2}\dot{4}$ \ansline[1.6cm]
%   \item $0.\dot{8}\dot{1}$ \ansline[1.6cm]
%   \item $0.\dot{6}$ \ansline[1.6cm]
%   \item $0.\dot{1}3\dot{5}$ \ansline[1.6cm]
%   \item $0.\dot{4}\dot{5}$ \ansline[1.6cm]
% \end{enumerate}
% \end{multicols}

% \end{enumerate}

% The nines rule is quick, but it only works when the repeating block starts
% immediately after the point --- it is no use at all on $0.1\dot{6}$ --- and it
% does not explain \emph{why} the nines appear. For both of those we need
% algebra, and that is the rest of the sheet.

\newpage

\parttitle{2}{Finding \ealkey{fractions} from \ealkey{recurring decimals}}

\ealpara{ترکیب یہ ہے کہ دو ایسے اعداد تلاش کریں جن کے تکراری دموں (tails) بالکل یکساں ہوں، اور پھر ایک کو دوسرے سے تفریق کریں۔ دم کٹ جاتے ہیں، اور جو بچتا ہے اس میں کوئی اعشاریے نہیں ہوتے۔}{The trick is to find two numbers whose recurring tails are \textbf{exactly the same}, and then subtract one from the other. The tails cancel, and what is left has no decimals in it at all.}
\ealpara{ایسا کرنے کے لیے ہمیں ایک \ealkeytr{متغیر} استعمال کرنا ہوگا، جسے ہم $x$ کہیں گے۔ $x$ اس \ealkeytr{تکراری اعشاریہ} کی نمائندگی کرتا ہے جسے ہم \ealkeytr{کسر} میں بدلنے کی کوشش کر رہے ہیں۔}{To do this we have to use an \textbf{algebraic \ealkey{variable}}, which we will call $x$. $x$ stands for the \ealkey{recurring decimal} we are trying to convert to a \ealkey{fraction}.}

\begin{wex}{turn $0.\dot{4}$ into a \ealkey{fraction}}
\small
\ealpara{\textbf{پہلا قدم --- $x$ متعارف کریں} ہم $x=0.\dot{4}$ لکھتے ہیں، تو $x=0.4444\ldots$}{\textbf{Step 1 --- introduce $x$} We write $x=0.\dot{4}$, so $x=0.4444\ldots$}

\smallskip
\ealpara{\textbf{دوسرا قدم --- مقامی قدریں سرکائیں۔} $10$ سے ضرب دینے سے ہر \ealkeytr{ہندسہ} ایک مقام بائیں طرف سرک جاتا ہے: \newline $10x=4.4444\ldots$ (یاد رکھیں $10x = 10 \times x$)}{\textbf{Step 2 --- shift the \ealkey{place values}.} Multiplying by $10$ moves every \ealkey{digit} one place to the left: \newline $10x=4.4444\ldots$ (remember $10x = 10 \times x$).}

\smallskip
\ealpara{\textbf{تیسرا قدم --- تفریق کریں۔} دونوں دم یکساں ہیں، اس لیے تفریق کرنے پر وہ کٹ جاتے ہیں: $10x -x$.}{\textbf{Step 3 --- subtract.} The two tails are identical, so they cancel out when we subtract: $10x -x$.}

\begin{center}
\begin{tikzpicture}[font=\small]
  \fill[yellow!45,rounded corners=2pt] \tailbox{0}{5}{-0.7}{0.26};
  \drow{0}{$10x\;=$}{$4$}{4,4,4,4,4,{$\ldots$}}
  \drow{-0.45}{$x\;=$}{$0$}{4,4,4,4,4,{$\ldots$}}
  \draw (-1.35,-0.72) -- (2.3,-0.72);
  \drow{-1.05}{$9x\;=$}{$4$}{0,0,0,0,0,{$\ldots$}}
  \node[anchor=west,font=\scriptsize,text width=6.2cm] at (2.7,-0.3)
    {the shaded tails are the same digits in the same columns, so they
     subtract to nothing};
\end{tikzpicture}
\end{center}
\vspace{-2pt}
\ealpara{تو $10x-x=4$، جس کا مطلب ہے $9x=4$۔}{So $10x-x=4$, which means $9x=4$.}

\smallskip
\ealpara{\textbf{چوتھا قدم --- حل کریں۔} ہمارے پاس $x$ کے $9$ حصے ہیں، اس لیے ہم دونوں اطراف کو $9$ سے تقسیم کرتے ہیں:}{\textbf{Step 4 --- solve.} We have $9$ lots of $x$, so we divide both sides by $9$:}

\vspace{-2pt}
\begin{center}
\divideboth{9x}{4}{x}{\dfrac{4}{9}}{9}
\end{center}
\end{wex}

\begin{enumerate}[resume]

\item \ealpara{خالی جگہیں بھریں تاکہ $0.\dot{7}$ کو \ealkeytr{کسر} میں بدلا جا سکے۔}{Fill in the gaps to turn $0.\dot{7}$ into a \ealkey{fraction}.}

\ealpara{\textbf{پہلا قدم --- $x$ متعارف کریں} ہم $x=0.\dot{7}$ لکھتے ہیں، تو $x=\ansline[3.5cm]$}{\textbf{Step 1 --- introduce $x$} We write $x=0.\dot{7}$, so $x=\ansline[3.5cm]$}

\smallskip
\ealpara{\textbf{دوسرا قدم --- مقامی قدریں سرکائیں۔} $10$ سے ضرب دینے سے ہر \ealkeytr{ہندسہ} ایک مقام بائیں طرف سرک جاتا ہے: \newline $10x=\ansline[3.5cm]$ (یاد رکھیں $10x = 10 \times x$)}{\textbf{Step 2 --- shift the \ealkey{place values}.} Multiplying by $10$ moves every \ealkey{digit} one place to the left: \newline $10x=\ansline[3.5cm]$ (remember $10x = 10 \times x$).}

\smallskip
\ealpara{\textbf{تیسرا قدم --- تفریق کریں۔} دونوں دم یکساں ہیں، اس لیے تفریق کرنے پر وہ کٹ جاتے ہیں: $10x-x$.}{\textbf{Step 3 --- subtract.} The two tails are identical, so they cancel out when we subtract: $10x-x$.}

\begin{center}
\begin{tikzpicture}[font=\small]
  \fill[yellow!45,rounded corners=2pt] \tailbox{0}{5}{-0.7}{0.26};
  \drowgap{0}{$10x\;=$}{\ansline[0.4cm]}{5}
  \drowgap{-0.45}{$x\;=$}{$0$}{5}
  \draw (-1.35,-0.72) -- (2.3,-0.72);
  \drowgap{-1.05}{$9x\;=$}{\ansline[0.4cm]}{5}
  \node[anchor=west,font=\scriptsize,text width=6.2cm] at (2.7,-0.3)
    {the shaded tails are the same digits in the same columns, so they
     subtract to nothing};
\end{tikzpicture}
\end{center}
\vspace{-2pt}
\ealpara{تو $10x-x=\ansline[1.6cm]$، جس کا مطلب ہے $9x=\ansline[1.6cm]$۔}{So $10x-x=\ansline[1.6cm]$, which means $9x=\ansline[1.6cm]$.}

\smallskip
\ealpara{\textbf{چوتھا قدم --- حل کریں۔} ہمارے پاس $x$ کے \ansline[0.9cm] حصے ہیں، اس لیے ہم دونوں اطراف کو \ansline[0.9cm] سے تقسیم کرتے ہیں۔ اپنا حل کام شدہ مثال کی طرح لکھیں۔}{\textbf{Step 4 --- solve.} We have \ansline[0.9cm] lots of $x$, so we divide both sides by \ansline[0.9cm]. Set your working out like the worked example.}

\vspace{13em}

\newpage

\item \ealpara{$0.\dot{8}$ کو \ealkeytr{کسر} میں بدلیں، وہی چار قدم استعمال کرتے ہوئے۔}{Turn $0.\dot{8}$ into a \ealkey{fraction}, using the same four steps.}
\begin{itemize}[label={},leftmargin=1.8em,itemsep=12pt,topsep=6pt]
  \item \ealpara{\textbf{پہلا قدم --- $x$ متعارف کریں} \hfill \ansline[7cm]}{\textbf{Step 1 --- introduce $x$} \hfill \ansline[7cm]}
  \item \ealpara{\textbf{دوسرا قدم --- مقامی قدریں سرکائیں۔} \hfill \ansline[7cm]}{\textbf{Step 2 --- shift the \ealkey{place values}.} \hfill \ansline[7cm]}
  \item \ealpara{\textbf{تیسرا قدم --- تفریق کریں۔} دونوں دم یکساں ہیں، اس لیے تفریق کرنے پر وہ کٹ جاتے ہیں: $10x-x$.}{\textbf{Step 3 --- subtract.} The two tails are identical, so they cancel out when we subtract: $10x-x$.}
        \begin{center}
        \begin{tikzpicture}[font=\small]
          \fill[yellow!45,rounded corners=2pt] \tailbox{0}{5}{-0.7}{0.26};
          \drowgap{0}{$10x\;=$}{\ansline[0.4cm]}{5}
          \drowgap{-0.45}{$x\;=$}{$0$}{5}
          \draw (-1.35,-0.72) -- (2.3,-0.72);
          \drowgap{-1.05}{$9x\;=$}{\ansline[0.4cm]}{5}
          \node[anchor=west,font=\scriptsize,text width=6.2cm] at (2.7,-0.3)
            {the shaded tails are the same digits in the same columns, so they
             subtract to nothing};
        \end{tikzpicture}
        \end{center}
        \vspace{-2pt}
        \ealpara{تو $10x-x=\ansline[1.6cm]$، جس کا مطلب ہے $9x=\ansline[1.6cm]$۔}{So $10x-x=\ansline[1.6cm]$, which means $9x=\ansline[1.6cm]$.}
  \item \ealpara{\textbf{چوتھا قدم --- حل کریں۔} دونوں اطراف کو \ansline[0.9cm] سے تقسیم کریں۔}{\textbf{Step 4 --- solve.} Divide both sides by \ansline[0.9cm].}

        \vspace{11em}
\end{itemize}

\item \ealpara{ان میں سے ہر ایک کو \ealkeytr{سادہ ترین صورت} میں \ealkeytr{کسر} میں بدلیں۔ ہر بار چاروں قدم لکھیں۔}{Turn each of these into a \ealkey{fraction} in its \ealkey{simplest form}. Write out all four steps each time.}
\begin{multicols}{4}
\begin{enumerate}
  \item $0.\dot{2}$
  \item $0.\dot{5}$
  \item $0.\dot{6}$
  \item $0.\dot{3}$
\end{enumerate}
\end{multicols}
\vspace{28em}

\end{enumerate}

\newpage

\parttitle{3}{Longer repeating blocks}

\ealpara{$0.\dot{3}\dot{6}$ کو اسی طرح آزمائیں اور دم کٹتے نہیں ہیں۔ $10$ سے ضرب دینے سے مقامی قدریں ایک جگہ سرک جاتی ہیں، جس سے گروہ کے دو \ealkeytr{ہندسے} غلط ترتیب میں آ جاتے ہیں۔ \emph{دو} \ealkeytr{ہندسوں} کے گروہ کو \emph{دو} مقامی قدروں سے سرکانا ہوتا ہے، اس لیے ہم اس کے بجائے $100$ سے ضرب دیتے ہیں۔}{Try $0.\dot{3}\dot{6}$ the same way and the tails do not cancel. Multiplying by $10$ shifts the \ealkey{place values} by one, which puts the two \ealkey{digits} of the block the wrong way round. A block of \emph{two} \ealkey{digits} has to be shifted by \emph{two} \ealkey{place values}, so we multiply by $100$ instead.}

\vspace{4pt}
\begin{center}
\begin{tikzpicture}[font=\small]
  \fill[red!18,rounded corners=2pt] \tailbox{0}{4}{-0.7}{0.26};
  \drow{0}{$10x\;=$}{$3$}{6,3,6,3,{$\ldots$}}
  \drow{-0.45}{$x\;=$}{$0$}{3,6,3,6,{$\ldots$}}
  \node[font=\scriptsize,anchor=north,red!60!black] at (0.4,-0.82)
    {\textbf{tails differ} --- nothing cancels};
  \begin{scope}[xshift=7.6cm]
    \fill[green!30,rounded corners=2pt] \tailbox{0}{4}{-0.7}{0.26};
    \drow{0}{$100x\;=$}{$36$}{3,6,3,6,{$\ldots$}}
    \drow{-0.45}{$x\;=$}{$0$}{3,6,3,6,{$\ldots$}}
    \node[font=\scriptsize,anchor=north,green!35!black] at (0.4,-0.82)
      {\textbf{tails match} --- these cancel};
  \end{scope}
\end{tikzpicture}
\end{center}
\vspace{4pt}

\begin{wex}{turn $0.\dot{3}\dot{6}$ into a \ealkey{fraction}}
\small
\ealpara{\textbf{پہلا قدم --- $x$ متعارف کریں} $x=0.\dot{3}\dot{6}$، تو $x=0.363636\ldots$ \\[8pt]
\textbf{دوسرا قدم --- مقامی قدریں سرکائیں۔} $100x=36.363636\ldots$ گروہ $2$ \ealkeytr{ہندسوں} کا ہے \\[8pt]
\textbf{تیسرا قدم --- تفریق کریں۔} $100x-x=36$ تو $99x=36$ \\[8pt]
\textbf{چوتھا قدم --- حل کریں۔} $x=\dfrac{36}{99}=\dfrac{4}{11}$}{%
\begin{tabular}{@{}ll@{\hspace{1.6em}}l@{}}
\textbf{Step 1 --- introduce $x$} & $x=0.\dot{3}\dot{6}$, so $x=0.363636\ldots$ & \\[8pt]
\textbf{Step 2 --- shift the \ealkey{place values}.} & $100x=36.363636\ldots$
  & the block is $2$ \ealkey{digits} long\\[8pt]
\textbf{Step 3 --- subtract.} & $100x-x=36$ & so $99x=36$\\[8pt]
\textbf{Step 4 --- solve.} & $x=\dfrac{36}{99}=\dfrac{4}{11}$ & \\
\end{tabular}
}
\end{wex}

\begin{enumerate}[resume]

\item \ealpara{خالی جگہیں بھریں تاکہ $0.\dot{7}\dot{2}$ کو \ealkeytr{کسر} میں بدلا جا سکے۔}{Fill in the gaps to turn $0.\dot{7}\dot{2}$ into a \ealkey{fraction}.}
\begin{itemize}[label={},leftmargin=1.8em,itemsep=12pt,topsep=6pt]
  \item \ealpara{\textbf{پہلا قدم --- $x$ متعارف کریں} \hfill $x=0.727272\ldots$}{\textbf{Step 1 --- introduce $x$} \hfill $x=0.727272\ldots$}
  \item \ealpara{\textbf{دوسرا قدم --- مقامی قدریں سرکائیں۔} گروہ \ansline[1cm] \ealkeytr{ہندسوں} کا ہے، تو $\ansline[1cm]x=\ansline[2.2cm]$}{\textbf{Step 2 --- shift the \ealkey{place values}.} The block is \ansline[1cm] \ealkey{digits} long, so $\ansline[1cm]x=\ansline[2.2cm]$}
  \item \ealpara{\textbf{تیسرا قدم --- تفریق کریں۔} \hfill $\ansline[1cm]x=\ansline[3cm]$}{\textbf{Step 3 --- subtract.} \hfill $\ansline[1cm]x=\ansline[3cm]$}
  \item \ealpara{\textbf{چوتھا قدم --- حل کریں۔} \hfill $x=\ansline[3cm]$}{\textbf{Step 4 --- solve.} \hfill $x=\ansline[3cm]$}
\end{itemize}

\item \ealpara{انہیں \ealkeytr{سادہ ترین صورت} میں \ealkeytr{کسروں} میں بدلیں۔ حصہ (e) میں تین \ealkeytr{ہندسوں} کا گروہ ہے، اس لیے سوچیں کہ کس سے ضرب دینا ہے۔}{Turn these into \ealkey{fractions} in their \ealkey{simplest form}. Part (e) has a block of three \ealkey{digits}, so think about what to multiply by.}
\begin{multicols}{5}
\begin{enumerate}
  \item $0.\dot{1}\dot{5}$
  \item $0.\dot{0}\dot{6}$
  \item $0.\dot{5}\dot{4}$
  \item $0.\dot{9}\dot{0}$
  \item $0.\dot{1}0\dot{8}$
\end{enumerate}
\end{multicols}
\vspace{13em}

\newpage

\item
\begin{enumerate}
  \item \ealpara{اپنے الفاظ میں سمجھائیں کہ آپ کیسے فیصلہ کرتے ہیں کہ $x$ کو کس سے ضرب دینا ہے۔}{Explain, in your own words, how you decide what to multiply $x$ by.}

        \vspace{4em}
  \item \ealpara{کیلا $0.\dot{1}\dot{2}$ کو تبدیل کرتی ہے اور $\frac{12}{99}$ حاصل کرتی ہے۔ بین وہی عدد تبدیل کرتا ہے اور $\frac{4}{33}$ حاصل کرتا ہے۔ کون صحیح ہے؟}{Kayla converts $0.\dot{1}\dot{2}$ and gets $\frac{12}{99}$. Ben converts the same number and gets $\frac{4}{33}$. Who is right?}

        \vspace{5em}
\end{enumerate}

\end{enumerate}

\parttitle{4}{When the repeating part does not start straight away}

\ealpara{$0.1\dot{6}$ زیادہ مشکل ہے، کیونکہ $1$ دہائیوں کے کالم میں ہے اور دہرایا نہیں جاتا۔ $10$ سے ضرب دینے سے ایسا عدد ملتا ہے جس کا دم تمام چھکے ہیں، لیکن وہ دم $x$ کے اپنے دم سے میل نہیں کھاتا۔ اس کے بجائے ہم مقامی قدروں کو \emph{دو بار} سرکاتے ہیں، اور دو نتائج کو ایک دوسرے سے تفریق کرتے ہیں۔}{$0.1\dot{6}$ is harder, because the $1$ is in the tenths column and does not repeat. Multiplying by $10$ gives a number whose tail is all sixes, but that tail does not match the tail of $x$ itself. Instead we shift the \ealkey{place values} \emph{twice}, and subtract the two results from each other.}

\vspace{2pt}
\begin{center}
\begin{tikzpicture}[font=\small]
  \fill[yellow!45,rounded corners=2pt] \tailbox{0}{5}{-1.15}{-0.24};
  \drow{0}{$x\;=$}{$0$}{1,6,6,6,6,{$\ldots$}}
  \drow{-0.45}{$10x\;=$}{$1$}{6,6,6,6,6,{$\ldots$}}
  \drow{-0.9}{$100x\;=$}{$16$}{6,6,6,6,6,{$\ldots$}}
  \node[anchor=west,font=\scriptsize,text width=5.4cm] at (2.6,0)
    {the tail of $x$ starts with a $1$, so it will never match};
  \node[anchor=west,font=\scriptsize,text width=5.4cm] at (2.6,-0.72)
    {\textbf{these two match}, so subtract these two};
\end{tikzpicture}
\end{center}

\begin{wex}{turn $0.1\dot{6}$ into a \ealkey{fraction}}
\small
\ealpara{\textbf{پہلا قدم --- $x$ متعارف کریں} $x=0.1\dot{6}$، تو $x=0.16666\ldots$ \\[8pt]
\textbf{دوسرا قدم --- مقامی قدریں سرکائیں۔} $10x=1.6666\ldots$ $1$ سے آگے سرکائیں \\[8pt]
$100x=16.6666\ldots$ ایک اور مقام سرکائیں \\[8pt]
\textbf{تیسرا قدم --- تفریق کریں۔} $100x-10x=16-1=15$ تو $90x=15$ \\[8pt]
\textbf{چوتھا قدم --- حل کریں۔} $x=\dfrac{15}{90}=\dfrac{1}{6}$}{%
\begin{tabular}{@{}ll@{\hspace{1.6em}}l@{}}
\textbf{Step 1 --- introduce $x$} & $x=0.1\dot{6}$, so $x=0.16666\ldots$ &\\[8pt]
\textbf{Step 2 --- shift the \ealkey{place values}.} & $10x=1.6666\ldots$
  & shift past the $1$\\[8pt]
 & $100x=16.6666\ldots$ & shift one more place\\[8pt]
\textbf{Step 3 --- subtract.} & $100x-10x=16-1=15$ & so $90x=15$\\[8pt]
\textbf{Step 4 --- solve.} & $x=\dfrac{15}{90}=\dfrac{1}{6}$ &\\
\end{tabular}
}
\end{wex}

\begin{enumerate}[resume]

\item \ealpara{خالی جگہیں بھریں تاکہ $0.2\dot{7}$ کو \ealkeytr{کسر} میں بدلا جا سکے۔}{Fill in the gaps to turn $0.2\dot{7}$ into a \ealkey{fraction}.}
\begin{itemize}[label={},leftmargin=1.8em,itemsep=12pt,topsep=6pt]
  \item \ealpara{\textbf{پہلا قدم --- $x$ متعارف کریں} \hfill $x=0.27777\ldots$}{\textbf{Step 1 --- introduce $x$} \hfill $x=0.27777\ldots$}
  \item \ealpara{\textbf{دوسرا قدم --- مقامی قدریں سرکائیں۔} $2$ سے آگے سرکائیں: \hfill $10x=\ansline[3cm]$}{\textbf{Step 2 --- shift the \ealkey{place values}.} Shift past the $2$: \hfill $10x=\ansline[3cm]$}
  \item \ealpara{\hspace*{1.2em}ایک اور مقام سرکائیں: \hfill $100x=\ansline[3cm]$}{\hspace*{1.2em}Shift one more place: \hfill $100x=\ansline[3cm]$}
  \item \ealpara{\textbf{تیسرا قدم --- تفریق کریں۔} \hfill $100x-10x=\ansline[3cm]$}{\textbf{Step 3 --- subtract.} \hfill $100x-10x=\ansline[3cm]$}
  \item \ealpara{\hspace*{1.2em}تو \hfill $90x=\ansline[3cm]$}{\hspace*{1.2em}so \hfill $90x=\ansline[3cm]$}
  \item \ealpara{\textbf{چوتھا قدم --- حل کریں۔} \hfill $x=\ansline[3cm]$}{\textbf{Step 4 --- solve.} \hfill $x=\ansline[3cm]$}
\end{itemize}

\newpage

\item \ealpara{انہیں \ealkeytr{سادہ ترین صورت} میں \ealkeytr{کسروں} میں بدلیں۔ (d) میں دہرایا جانے والا گروہ دو \ealkeytr{ہندسوں} کا ہے، اس لیے غور سے سوچیں کہ $x$ کے کون سے دو ضرب استعمال کرنے ہیں۔}{Turn these into \ealkey{fractions} in their \ealkey{simplest form}. In (d) the repeating block is two \ealkey{digits} long, so think carefully about which two multiples of $x$ you need.}
\begin{multicols}{4}
\begin{enumerate}
  \item $0.4\dot{3}$
  \item $0.58\dot{3}$
  \item $0.08\dot{3}$
  \item $0.1\dot{2}\dot{4}$
\end{enumerate}
\end{multicols}
\vspace{13em}

\item \ealpara{وہی چار قدم اس وقت بھی کام کرتے ہیں جب عدد $1$ سے بڑا ہو۔ انہیں \ealkeytr{سادہ ترین صورت} میں \ealkeytr{غیر مناسب کسروں} میں بدلیں۔}{The same four steps work when the number is bigger than $1$. Turn these into \ealkey{improper fractions} in their \ealkey{simplest form}.}
\begin{multicols}{3}
\begin{enumerate}
  \item $2.\dot{7}$
  \item $1.\dot{4}\dot{5}$
  \item $3.\dot{1}\dot{8}$
\end{enumerate}
\end{multicols}
\vspace{13em}

\item \ealpara{\textbf{مخلوط مشق۔} خود فیصلہ کریں کہ ہر بار کس سے ضرب دینا ہے۔ ہر جواب کو سادہ کریں۔}{\textbf{Mixed practice.} Decide for yourself what to multiply by each time. Simplify every answer.}
\begin{multicols}{4}
\begin{enumerate}
  \item $0.\dot{8}$
  \item $0.\dot{6}\dot{3}$
  \item $0.2\dot{5}$
  \item $0.\dot{0}2\dot{7}$
  \item $0.7\dot{2}$
  \item $0.\dot{1}\dot{2}$
  \item $0.41\dot{6}$
  \item $0.06\dot{4}$
\end{enumerate}
\end{multicols}
\vspace{13em}

\end{enumerate}

\newpage

\parttitle{5}{Challenges}

\begin{enumerate}[resume]

\item \ealpara{\textbf{$0.\dot{9}$ کی قدر۔} حصہ 2 کے چار قدم $x=0.\dot{9}$ پر استعمال کریں۔ آپ کو کون سی \ealkeytr{کسر} ملتی ہے؟}{\textbf{The value of $0.\dot{9}$.} Use the four steps from Part 2 on $x=0.\dot{9}$. What \ealkey{fraction} do you get?}
\begin{enumerate}
  \item 
        \vspace{12em}
  \item \ealpara{ایلا کہتی ہے: ''یہ غلط ہونا چاہیے۔ $0.9999\ldots$ ظاہر ہے $1$ سے تھوڑا کم ہے۔'' نیچے دی گئی عددی لکیر $0.9$، پھر $0.99$، پھر $0.999$، اور اسی طرح نشان زد کرتی ہے۔ اسے استعمال کریں، اس حقیقت کے ساتھ کہ $\frac{1}{3}=0.\dot{3}$، اور ایلا کو جواب لکھیں۔}{Ella says: ``That has to be wrong. $0.9999\ldots$ is obviously a bit less than $1$.'' The number line below marks $0.9$, then $0.99$, then $0.999$, and so on. Use it, together with the fact that $\frac{1}{3}=0.\dot{3}$, to write a reply to Ella.}
\end{enumerate}

\vspace{4pt}
\begin{center}
\begin{tikzpicture}[font=\small]
  \node[font=\scriptsize] at (4.4,0.7) {each step along the line closes half of the gap that is left};
  \draw[-{Stealth[length=3mm]}] (-0.5,0) -- (10.2,0);
  \foreach \x/\lab in {0/$0.9$, 4.5/$0.99$, 6.75/$0.999$, 7.875/$0.9999$}{
    \draw (\x,0.18) -- (\x,-0.18);
    \node[below,font=\scriptsize] at (\x,-0.2) {\lab};}
  \foreach \x in {8.44,8.72,8.86,8.93}{\draw (\x,0.12) -- (\x,-0.12);}
  \draw[thick] (9,0.34) -- (9,-0.34);
  \node[below,font=\small] at (9,-0.36) {$1$};
\end{tikzpicture}
\end{center}
\vspace{4pt}

\begin{enumerate}[resume]
  \item \ealpara{سب سے بڑا عدد کون سا ہے جو $1$ سے چھوٹا ہو؟ سمجھائیں۔}{What is the largest number that is smaller than $1$? Explain.}

        \vspace{14em}
\end{enumerate}

\newpage

\item \ealpara{\textbf{کون سی \ealkeytr{کسریں} دہرائی جاتی ہیں؟} کچھ \ealkeytr{کسریں} ایسے اعشاریے دیتی ہیں جو رک جاتے ہیں --- انہیں \emph{\ealkeytr{ختم ہونے والے اعشاریے}} کہا جاتا ہے --- اور کچھ ایسے اعشاریے دیتی ہیں جو ہمیشہ دہراتے رہتے ہیں۔}{\textbf{Which \ealkey{fractions} recur?} Some \ealkey{fractions} give decimals that stop --- these are called \emph{\ealkey{terminating decimals}} --- and some give decimals that recur forever.}
\begin{enumerate}
  \item \ealpara{مختصر تقسیم استعمال کرتے ہوئے جدول مکمل کریں۔ ختم ہونے والوں پر دائرہ لگائیں۔}{Complete the table using short division. Circle the terminating ones.}

\vspace{3pt}
\begin{center}\small
\begin{tabular}{|c|c|c|c|c|c|}
\hline
\rule{0pt}{2.8em}$\frac{1}{2}=0.5$ & $\frac{1}{3}=$\rule{1.48cm}{0pt} &
$\frac{1}{4}=$\rule{1.48cm}{0pt} & $\frac{1}{5}=$\rule{1.48cm}{0pt} &
$\frac{1}{6}=$\rule{1.48cm}{0pt} & $\frac{1}{7}=$\rule{1.48cm}{0pt}\\
\hline
\rule{0pt}{2.8em}$\frac{1}{8}=$\rule{1.48cm}{0pt} & $\frac{1}{9}=$\rule{1.48cm}{0pt} &
$\frac{1}{10}=$\rule{1.48cm}{0pt} & $\frac{1}{11}=$\rule{1.48cm}{0pt} &
$\frac{1}{12}=$\rule{1.48cm}{0pt} & $\frac{1}{16}=$\rule{1.48cm}{0pt}\\
\hline
\end{tabular}
\end{center}
\vspace{3pt}

  \item \ealpara{ہر ختم ہونے والے \ealkeytr{مخرج} کو اس کے \ealkeytr{مفرد عوامل} کے \ealkeytr{حاصل ضرب} کے طور پر لکھیں۔ ان سب میں کیا مشترک ہے؟}{Write each terminating \ealkey{denominator} as a \ealkey{product} of its \ealkey{prime factors}. What do they all have in common?}

        \vspace{9em}
  \item \ealpara{تقسیم کیے بغیر اندازہ لگائیں کہ ان میں سے ہر ایک ختم ہوتا ہے یا دہرایا جاتا ہے، اور بتائیں کہ آپ کیسے جانتے ہیں۔}{Without dividing, predict whether each of these terminates or recurs, and say how you know.}
        \[\frac{1}{15}\qquad\frac{1}{40}\qquad\frac{1}{35}\qquad\frac{1}{64}\]

        \vspace{9em}
  \item \ealpara{\ealkeytr{مخرج} میں صرف $2$s اور $5$s ہونے سے اعشاریہ کیوں رک جاتا ہے؟ \emph{اشارہ: $2\times5$ میں کیا خاص ہے؟}}{Why does having only $2$s and $5$s in the \ealkey{denominator} make the decimal stop? \emph{Hint: what is special about $2\times5$?}}

        \vspace{9em}
\end{enumerate}

\newpage

\item \ealpara{\textbf{ساتواں حصہ۔} ایک ساتواں حصہ چھ \ealkeytr{ہندسوں} کا دہرایا جانے والا گروہ رکھتا ہے: $\dfrac{1}{7}=0.\dot{1}4285\dot{7}$۔}{\textbf{Sevenths.} One seventh has a repeating block of six \ealkey{digits}: $\dfrac{1}{7}=0.\dot{1}4285\dot{7}$.}

\vspace{2pt}
\begin{minipage}[t]{0.58\linewidth}
\begin{enumerate}
  \item \ealpara{مختصر تقسیم استعمال کرتے ہوئے $\frac{2}{7}$، $\frac{3}{7}$، $\frac{4}{7}$، $\frac{5}{7}$ اور $\frac{6}{7}$ کو \ealkeytr{تکراری اعشاریوں} کے طور پر معلوم کریں۔}{Use short division to work out $\frac{2}{7}$, $\frac{3}{7}$, $\frac{4}{7}$, $\frac{5}{7}$ and $\frac{6}{7}$ as \ealkey{recurring decimals}.}

        \vspace{9em}
  \item \ealpara{آپ نے کیا محسوس کیا؟ جو ہوتا ہے اسے بیان کرنے کے لیے پہیے کو استعمال کریں۔}{What do you notice? Use the wheel to describe what happens.}

        \vspace{6em}
  \item \ealpara{$0.\dot{2}8571\dot{4}$ پر چار قدم استعمال کرتے ہوئے اسے واپس \ealkeytr{کسر} میں بدلیں، اور جانچیں کہ یہ حصہ (a) کے آپ کے جواب سے مطابقت رکھتا ہے۔}{Use the four steps on $0.\dot{2}8571\dot{4}$ to turn it back into a \ealkey{fraction}, and check that it agrees with your answer to part (a).}

        \vspace{7em}
\end{enumerate}
\end{minipage}\hfill
\begin{minipage}[t]{0.38\linewidth}
\centering
\begin{tikzpicture}[baseline=(current bounding box.north),font=\small]
  \foreach \d [count=\i from 0] in {1,4,2,8,5,7}{
    \node[circle,draw,minimum size=8mm,inner sep=0pt,fill=blue!8]
      (n\i) at ({90-\i*60}:1.6) {\d};}
  \foreach \i/\j in {0/1,1/2,2/3,3/4,4/5,5/0}{
    \draw[-{Stealth[length=2mm]},blue!45!black] (n\i) to[bend left=18] (n\j);}
  \node[font=\scriptsize,text width=4cm,align=center] at (0,-2.5)
    {reading round the wheel from a different starting digit gives a
     different seventh};
\end{tikzpicture}
\end{minipage}

\newpage

\item \ealpara{\textbf{اپنا بنائیں۔} ایک \ealkeytr{تکراری اعشاریہ} بنائیں جو \ealkeytr{مخرج} $11$ والی \ealkeytr{کسر} میں تبدیل ہو۔ دکھائیں کہ یہ کام کرتا ہے۔}{\textbf{Make your own.} Invent a \ealkey{recurring decimal} that converts to a \ealkey{fraction} with \ealkey{denominator} $11$. Show that it works.}
\begin{enumerate}
  \item 
        \vspace{10em}
  \item \ealpara{وہ \ealkeytr{تکراری اعشاریہ} تلاش کریں جو بالکل $\frac{5}{6}$ ہو، پھر چار قدم استعمال کرتے ہوئے اسے واپس $\frac{5}{6}$ میں بدلیں۔}{Find the \ealkey{recurring decimal} that is exactly $\frac{5}{6}$, then use the four steps to turn it back into $\frac{5}{6}$.}

        \vspace{10em}
  \item \ealpara{کیا \emph{ہر} \ealkeytr{تکراری اعشاریہ} کسی \ealkeytr{کسر} کے برابر ہوتا ہے؟ اپنے الفاظ میں سمجھائیں کہ تفریق ہمیشہ کیوں کام کرتی ہے، چاہے آپ کسی بھی \ealkeytr{تکراری اعشاریہ} سے شروع کریں۔}{Is \emph{every} \ealkey{recurring decimal} equal to a \ealkey{fraction}? Explain, in your own words, why the subtraction always works, whatever \ealkey{recurring decimal} you start with.}

        \vspace{10em}
  \item \ealpara{بین کہتا ہے: ''$\pi=3.14159\ldots$ ہمیشہ چلتا رہتا ہے، اس لیے کوئی \ealkeytr{تکراری اعشاریہ} ضرور ہوگا جو بالکل $\pi$ کے برابر ہو۔'' حصہ (c) کے اپنے جواب کا استعمال کرتے ہوئے سمجھائیں کہ بین غلط کیوں ہے۔}{Ben says: ``$\pi=3.14159\ldots$ goes on forever, so there must be a \ealkey{recurring decimal} that is exactly equal to $\pi$.'' Use your answer to part (c) to explain why Ben is wrong.}

        \vspace{10em}
\end{enumerate}

\end{enumerate}

\newpage

\heading{Answers}

{\small
\begin{multicols}{2}
\raggedcolumns
\setlength{\parskip}{4pt}

\ealpara{\textbf{سوال 1۔} (a) $0.55555555$ \quad (b) $0.18181818$ \quad (c) $0.37777777$ \quad (d) $0.20620620$ \quad (e) $0.45151515$ \quad (f) $0.99999999$}{\textbf{Q1.} (a) $0.55555555$ \quad (b) $0.18181818$ \quad (c) $0.37777777$ \quad (d) $0.20620620$ \quad (e) $0.45151515$ \quad (f) $0.99999999$}

\ealpara{\textbf{سوال 2۔} (a) $0.\dot{8}$ \quad (b) $0.\dot{6}\dot{3}$ \quad (c) $0.2\dot{4}$ \quad (d) $0.\dot{1}5\dot{7}$ \quad (e) $0.91\dot{6}$ \quad (f) $0.4\dot{0}\dot{5}$}{\textbf{Q2.} (a) $0.\dot{8}$ \quad (b) $0.\dot{6}\dot{3}$ \quad (c) $0.2\dot{4}$ \quad (d) $0.\dot{1}5\dot{7}$ \quad (e) $0.91\dot{6}$ \quad (f) $0.4\dot{0}\dot{5}$}

\ealpara{\textbf{سوال 3۔} (a) $0.\dot{4}\dot{5}=0.454545$ اور $0.4\dot{5}=0.455555$، اس لیے $0.4\dot{5}$ بڑا ہے --- وہ دو \ealkeytr{اعشاریہ مقامات} تک متفق ہیں، لیکن تیسرا \ealkeytr{ہندسہ} $5$ بمقابلہ $4$ ہے۔ \quad (b) بہت سے جوابات، مثلاً $0.\dot{3}\dot{5}=0.353535\ldots$}{\textbf{Q3.} (a) $0.\dot{4}\dot{5}=0.454545$ and $0.4\dot{5}=0.455555$, so $0.4\dot{5}$ is bigger --- they agree to two \ealkey{decimal places}, but the third \ealkey{digit} is $5$ against $4$. \quad (b) Many answers, for example $0.\dot{3}\dot{5}=0.353535\ldots$}

\ealpara{\textbf{سوال 4۔} $10x=7.7777\ldots$; \; $10x-x=7$; \; $9x=7$; \; $x=\tfrac{7}{9}$}{\textbf{Q4.} $10x=7.7777\ldots$; \; $10x-x=7$; \; $9x=7$; \; $x=\tfrac{7}{9}$}

\ealpara{\textbf{سوال 5۔} $x=0.8888\ldots$; \; $10x=8.8888\ldots$; \; $10x-x=8$, تو $9x=8$; \; $x=\tfrac{8}{9}$}{\textbf{Q5.} $x=0.8888\ldots$; \; $10x=8.8888\ldots$; \; $10x-x=8$, so $9x=8$; \; $x=\tfrac{8}{9}$}

\ealpara{\textbf{سوال 6۔} (a) $\tfrac{2}{9}$ \quad (b) $\tfrac{5}{9}$ \quad (c) $\tfrac{6}{9}=\tfrac{2}{3}$ \quad (d) $\tfrac{3}{9}=\tfrac{1}{3}$}{\textbf{Q6.} (a) $\tfrac{2}{9}$ \quad (b) $\tfrac{5}{9}$ \quad (c) $\tfrac{6}{9}=\tfrac{2}{3}$ \quad (d) $\tfrac{3}{9}=\tfrac{1}{3}$}

\ealpara{\textbf{سوال 7۔} گروہ $2$ \ealkeytr{ہندسوں} کا ہے، اس لیے $100$ سے ضرب دیں: $100x=72.727272\ldots$; \; $99x=72$; \; $x=\tfrac{72}{99}=\tfrac{8}{11}$}{\textbf{Q7.} The block is $2$ \ealkey{digits} long, so multiply by $100$: $100x=72.727272\ldots$; \; $99x=72$; \; $x=\tfrac{72}{99}=\tfrac{8}{11}$}

\ealpara{\textbf{سوال 8۔} (a) $\tfrac{15}{99}=\tfrac{5}{33}$ \quad (b) $\tfrac{6}{99}=\tfrac{2}{33}$ \quad (c) $\tfrac{54}{99}=\tfrac{6}{11}$ \quad (d) $\tfrac{90}{99}=\tfrac{10}{11}$ \quad (e) $\tfrac{108}{999}=\tfrac{4}{37}$، $1000$ سے ضرب دینا}{\textbf{Q8.} (a) $\tfrac{15}{99}=\tfrac{5}{33}$ \quad (b) $\tfrac{6}{99}=\tfrac{2}{33}$ \quad (c) $\tfrac{54}{99}=\tfrac{6}{11}$ \quad (d) $\tfrac{90}{99}=\tfrac{10}{11}$ \quad (e) $\tfrac{108}{999}=\tfrac{4}{37}$, multiplying by $1000$}

\ealpara{\textbf{سوال 9۔} (a) دہرائے جانے والے گروہ میں \ealkeytr{ہندسے} گنیں۔ اگر $n$ ہیں، تو $1$ کے بعد $n$ صفر سے ضرب دیں، تاکہ گروہ پوری بار سرک جائے اور دم دوبارہ سیدھ میں آ جائیں۔ \quad (b) دونوں صحیح ہیں: $\tfrac{12}{99}=\tfrac{4}{33}$۔ بین نے اسے سادہ کیا ہے، جو سوال مانگتا ہے۔}{\textbf{Q9.} (a) Count the \ealkey{digits} in the repeating block. If there are $n$ of them, multiply by $1$ followed by $n$ zeros, so that the block shifts a whole number of times and the tails line up again. \quad (b) Both are right: $\tfrac{12}{99}=\tfrac{4}{33}$. Ben has simplified his, which is what the question asks for.}

\ealpara{\textbf{سوال 10۔} $10x=2.7777\ldots$; \; $100x=27.7777\ldots$; \; $100x-10x=27-2=25$; \; $90x=25$; \; $x=\tfrac{25}{90}=\tfrac{5}{18}$}{\textbf{Q10.} $10x=2.7777\ldots$; \; $100x=27.7777\ldots$; \; $100x-10x=27-2=25$; \; $90x=25$; \; $x=\tfrac{25}{90}=\tfrac{5}{18}$}

\ealpara{\textbf{سوال 11۔} (a) $\tfrac{39}{90}=\tfrac{13}{30}$ \quad (b) $\tfrac{525}{900}=\tfrac{7}{12}$ \quad (c) $\tfrac{75}{900}=\tfrac{1}{12}$ \quad (d) $\tfrac{123}{990}=\tfrac{41}{330}$، $1000x-10x$ استعمال کرتے ہوئے}{\textbf{Q11.} (a) $\tfrac{39}{90}=\tfrac{13}{30}$ \quad (b) $\tfrac{525}{900}=\tfrac{7}{12}$ \quad (c) $\tfrac{75}{900}=\tfrac{1}{12}$ \quad (d) $\tfrac{123}{990}=\tfrac{41}{330}$, using $1000x-10x$}

\ealpara{\textbf{سوال 12۔} (a) $\tfrac{25}{9}$ \quad (b) $\tfrac{144}{99}=\tfrac{16}{11}$ \quad (c) $\tfrac{315}{99}=\tfrac{35}{11}$}{\textbf{Q12.} (a) $\tfrac{25}{9}$ \quad (b) $\tfrac{144}{99}=\tfrac{16}{11}$ \quad (c) $\tfrac{315}{99}=\tfrac{35}{11}$}

\ealpara{\textbf{سوال 13۔} (a) $\tfrac{8}{9}$ \quad (b) $\tfrac{63}{99}=\tfrac{7}{11}$ \quad (c) $\tfrac{23}{90}$ \quad (d) $\tfrac{27}{999}=\tfrac{1}{37}$ \quad (e) $\tfrac{65}{90}=\tfrac{13}{18}$ \quad (f) $\tfrac{12}{99}=\tfrac{4}{33}$ \quad (g) $\tfrac{375}{900}=\tfrac{5}{12}$ \quad (h) $\tfrac{58}{900}=\tfrac{29}{450}$}{\textbf{Q13.} (a) $\tfrac{8}{9}$ \quad (b) $\tfrac{63}{99}=\tfrac{7}{11}$ \quad (c) $\tfrac{23}{90}$ \quad (d) $\tfrac{27}{999}=\tfrac{1}{37}$ \quad (e) $\tfrac{65}{90}=\tfrac{13}{18}$ \quad (f) $\tfrac{12}{99}=\tfrac{4}{33}$ \quad (g) $\tfrac{375}{900}=\tfrac{5}{12}$ \quad (h) $\tfrac{58}{900}=\tfrac{29}{450}$}

\ealpara{\textbf{سوال 14۔} (a) $10x=9.9999\ldots$، تو $9x=9$ اور $x=\tfrac{9}{9}=1$۔ \quad (b) طریقہ درست ہے، اس لیے $0.\dot{9}$ واقعی $1$ ہے۔ عددی لکیر پر ہر قدم جو بچا ہے اس کا آدھا بند کر دیتا ہے، اس لیے نشان $1$ کے جتنے قریب چاہیں پہنچ جاتے ہیں --- $0.\dot{9}$ اور $1$ کے درمیان ان کے مختلف اعداد ہونے کی کوئی گنجائش نہیں بچتی۔ نیز $\frac{1}{3}=0.\dot{3}$، اور دونوں اطراف کو $3$ سے ضرب دینے پر $1=0.\dot{9}$ ملتا ہے۔ \quad (c) کوئی نہیں ہے۔ آپ $1$ سے نیچے جو بھی عدد بتائیں، اس کے اور $1$ کے درمیان ہمیشہ ایک اور عدد ہوتا ہے --- یہی وجہ ہے کہ $0.\dot{9}$ کو $1$ کے بالکل نیچے ہونے کے بجائے $1$ ہونا \emph{ہی} پڑتا ہے۔}{\textbf{Q14.} (a) $10x=9.9999\ldots$, so $9x=9$ and $x=\tfrac{9}{9}=1$. \quad (b) The method is sound, so $0.\dot{9}$ really is $1$. Every step along the number line closes half of what is left, so the marks get as close to $1$ as you like --- there is no room left between $0.\dot{9}$ and $1$ for them to be different numbers. Also $\frac{1}{3}=0.\dot{3}$, and multiplying both sides by $3$ gives $1=0.\dot{9}$. \quad (c) There isn't one. Whatever number you name below $1$, there is always another number between it and $1$ --- which is exactly why $0.\dot{9}$ has to \emph{be} $1$ rather than sit just below it.}

\ealpara{\textbf{سوال 15۔} (a) $\tfrac{1}{3}=0.\dot{3}$, $\tfrac{1}{4}=0.25$, $\tfrac{1}{5}=0.2$, $\tfrac{1}{6}=0.1\dot{6}$, $\tfrac{1}{7}=0.\dot{1}4285\dot{7}$, $\tfrac{1}{8}=0.125$, $\tfrac{1}{9}=0.\dot{1}$, $\tfrac{1}{10}=0.1$, $\tfrac{1}{11}=0.\dot{0}\dot{9}$, $\tfrac{1}{12}=0.08\dot{3}$, $\tfrac{1}{16}=0.0625$۔ ختم ہونے والے یہ ہیں: $\tfrac{1}{2},\tfrac{1}{4},\tfrac{1}{5},\tfrac{1}{8},\tfrac{1}{10},\tfrac{1}{16}$۔ \quad (b) $2$, $2^2$, $5$, $2^3$, $2\times5$, $2^4$ --- ہر ایک صرف $2$s اور $5$s سے بنا ہے۔ \quad (c) $\tfrac{1}{15}$ دہرایا جاتا ہے، کیونکہ $15=3\times5$ اور $3$ کی اجازت نہیں؛ $\tfrac{1}{40}$ ختم ہوتا ہے، کیونکہ $40=2^3\times5$; $\tfrac{1}{35}$ دہرایا جاتا ہے، کیونکہ $35=5\times7$; $\tfrac{1}{64}$ ختم ہوتا ہے، کیونکہ $64=2^6$۔ \quad (d) ختم ہونے والا اعشاریہ صرف $10$، $100$، $1000$، \ldots کے اوپر ایک \ealkeytr{کسر} ہوتا ہے، اور $10=2\times5$۔ اس لیے ایک \ealkeytr{کسر} کو دس کی قوت کے اوپر دوبارہ لکھا جا سکتا ہے جب اس کا \ealkeytr{مخرج} $2$s اور $5$s سے بنا ہو؛ کوئی اور مفرد عامل کبھی کٹ نہیں سکتا۔}{\textbf{Q15.} (a) $\tfrac{1}{3}=0.\dot{3}$, $\tfrac{1}{4}=0.25$, $\tfrac{1}{5}=0.2$, $\tfrac{1}{6}=0.1\dot{6}$, $\tfrac{1}{7}=0.\dot{1}4285\dot{7}$, $\tfrac{1}{8}=0.125$, $\tfrac{1}{9}=0.\dot{1}$, $\tfrac{1}{10}=0.1$, $\tfrac{1}{11}=0.\dot{0}\dot{9}$, $\tfrac{1}{12}=0.08\dot{3}$, $\tfrac{1}{16}=0.0625$. The terminating ones are $\tfrac{1}{2},\tfrac{1}{4},\tfrac{1}{5},\tfrac{1}{8},\tfrac{1}{10},\tfrac{1}{16}$. \quad (b) $2$, $2^2$, $5$, $2^3$, $2\times5$, $2^4$ --- every one is built out of $2$s and $5$s only. \quad (c) $\tfrac{1}{15}$ recurs, since $15=3\times5$ and the $3$ is not allowed; $\tfrac{1}{40}$ terminates, since $40=2^3\times5$; $\tfrac{1}{35}$ recurs, since $35=5\times7$; $\tfrac{1}{64}$ terminates, since $64=2^6$. \quad (d) A terminating decimal is just a \ealkey{fraction} over $10$, $100$, $1000$, \ldots, and $10=2\times5$. So a \ealkey{fraction} can be rewritten over a \ealkey{power} of ten exactly when its \ealkey{denominator} is built from $2$s and $5$s; any other \ealkey{prime factor} can never be cancelled away.}

\ealpara{\textbf{سوال 16۔} (a) $\tfrac{2}{7}=0.\dot{2}8571\dot{4}$, $\tfrac{3}{7}=0.\dot{4}2857\dot{1}$, $\tfrac{4}{7}=0.\dot{5}7142\dot{8}$, $\tfrac{5}{7}=0.\dot{7}1428\dot{5}$, $\tfrac{6}{7}=0.\dot{8}5714\dot{2}$۔ \quad (b) وہ سب وہی چھ \ealkeytr{ہندسے} اسی چکری ترتیب میں استعمال کرتے ہیں، $1\to4\to2\to8\to5\to7\to1$؛ ہر ساتواں حصہ پہیے پر ایک مختلف جگہ سے شروع ہوتا ہے۔ \quad (c) گروہ $6$ \ealkeytr{ہندسوں} کا ہے، اس لیے $1\,000\,000$ سے ضرب دیں: اس سے $999\,999x=285\,714$ ملتا ہے، تو $x=\tfrac{285714}{999999}=\tfrac{2}{7}$، اوپر اور نیچے $142\,857$ سے تقسیم کرتے ہوئے۔}{\textbf{Q16.} (a) $\tfrac{2}{7}=0.\dot{2}8571\dot{4}$, $\tfrac{3}{7}=0.\dot{4}2857\dot{1}$, $\tfrac{4}{7}=0.\dot{5}7142\dot{8}$, $\tfrac{5}{7}=0.\dot{7}1428\dot{5}$, $\tfrac{6}{7}=0.\dot{8}5714\dot{2}$. \quad (b) They all use the same six \ealkey{digits} in the same cyclic order, $1\to4\to2\to8\to5\to7\to1$; each seventh simply starts at a different place on the wheel. \quad (c) The block is $6$ \ealkey{digits} long, so multiply by $1\,000\,000$: this gives $999\,999x=285\,714$, so $x=\tfrac{285714}{999999}=\tfrac{2}{7}$, dividing top and bottom by $142\,857$.}

\ealpara{\textbf{سوال 17۔} (a) بہت سے جوابات۔ مثلاً $0.\dot{2}\dot{7}$ سے $99x=27$ ملتا ہے، تو $x=\tfrac{27}{99}=\tfrac{3}{11}$؛ کوئی بھی دو \ealkeytr{ہندسوں} کا گروہ جو $9$ کا ضرب ہو لیکن $11$ کا نہ ہو، چلے گا۔ \quad (b) $\tfrac{5}{6}=0.8\dot{3}$، اور جانچ: $100x-10x=83-8=75$، تو $90x=75$ اور $x=\tfrac{75}{90}=\tfrac{5}{6}$۔ \quad (c) ہاں۔ اعشاریہ جو بھی ہو، اسے دہرائے جانے والے گروہ کی لمبائی سے سرکانے پر ہمیشہ ایک دوسرا عدد ملتا ہے جس کا دم یکساں ہوتا ہے، اس لیے تفریق ہمیشہ ایک پورا عدد چھوڑتی ہے۔ اس سے (پورا عدد)$\,\times x=$ (پورا عدد) کی شکل کی \ealkeytr{مساوات} ملتی ہے، اور تقسیم کرنے پر \ealkeytr{کسر} ملتی ہے۔ \quad (d) $\pi$ کوئی \ealkeytr{کسر} نہیں ہے، اور حصہ (c) ظاہر کرتا ہے کہ ہر \ealkeytr{تکراری اعشاریہ} ایک \ealkeytr{کسر} ہوتا ہے۔ اس لیے کوئی \ealkeytr{تکراری اعشاریہ} $\pi$ کے برابر نہیں ہو سکتا: اس کے \ealkeytr{ہندسے} ہمیشہ چلتے رہتے ہیں اور کبھی دہرانے والے گروہ میں نہیں آتے۔}{\textbf{Q17.} (a) Many answers. For example $0.\dot{2}\dot{7}$ gives $99x=27$, so $x=\tfrac{27}{99}=\tfrac{3}{11}$; any two-\ealkey{digit} block that is a multiple of $9$ but not of $11$ will do. \quad (b) $\tfrac{5}{6}=0.8\dot{3}$, and checking: $100x-10x=83-8=75$, so $90x=75$ and $x=\tfrac{75}{90}=\tfrac{5}{6}$. \quad (c) Yes. Whatever the decimal, shifting it by the length of the repeating block always produces a second number with an identical tail, so the subtraction always leaves a whole number. That gives an \ealkey{equation} of the form (whole number)$\,\times x=$ (whole number), and dividing gives a \ealkey{fraction}. \quad (d) $\pi$ is not a \ealkey{fraction}, and part (c) shows that every \ealkey{recurring decimal} is one. So no \ealkey{recurring decimal} can equal $\pi$: its \ealkey{digits} go on forever without ever settling into a repeating block.}

\end{multicols}
}

\end{document}